\chapter{Approssimanti di Padé}
\section{Introduzione}
Gli \textbf{approssimanti di Padé} rappresentano la migliore approssimazione di una funzione $f(\theta)$  ottenibile tramite una funzione razionale di ordine assegnato. Essi sono definiti dalla relazione:

\[R_{(S, T)}(\theta) = \frac{\sum_{i=0}^{T} p_i \theta^i}{1 + \sum_{k=1}^{S} q_k \theta^k} = \frac{p_0 + p_1 \theta + p_2 \theta^2 + \cdots + p_T \theta^T}{1 + q_1 \theta + q_2 \theta^2 + \cdots + q_S \theta^S}\ \ .\]

I coefficienti  $p_i$  e  $q_k$  sono determinati in modo da soddisfare, al massimo ordine possibile, le seguenti identità:
\[\left. \frac{\mathrm{d}^n R}{\mathrm{d} \theta^n} \right|_0 = \left. \frac{\mathrm{d}^n f}{\mathrm{d} \theta^n} \right|_0 \quad (n = 0, \dots, S + T)\]

%
%
%  
\section{Applicazione alla funzione esponenziale}

\subsection{Approssimante $ R_{(0,1)}(\theta)$ }

Consideriamo la funzione esponenziale  $e^{\theta}$  e definiamo l’approssimante  $R_{(0,1)}(\theta)$  come:

\[R_{(0,1)}(\theta) = p_0 + p_1 \theta  \]

Imponendo le condizioni di approssimazione di cui sopra:

\[R_{(0,1)}(0) = \left. e^{\theta} \right|_0 \implies p_0 = 1  \]

\[\left. \frac{\mathrm{d} R_{(0,1)}}{\mathrm{d} \theta} \right|_0 = \left. \frac{\mathrm{d} e^{\theta}}{\mathrm{d} \theta} \right|_0 \implies p_1 = 1  \]

Si ottiene quindi:

\[R_{(0,1)}(\theta) = 1 + \theta  \]

%
\subsection{Approssimante  $R_{(1,0)}(\theta)$ }
Definiamo ora l’approssimante  $R_{(1,0)}(\theta)$  come:

\[R_{(1,0)}(\theta) = \frac{p_0}{1 + q_1 \theta}  \]

Imponendo le condizioni di approssimazione:

\[R_{(1,0)}(0) = \left. e^{\theta} \right|_0 \implies p_0 = 1  \]

\[\left. \frac{\mathrm{d} R_{(1,0)}}{\mathrm{d} \theta} \right|_0 = \left. 
\frac{\mathrm{d} e^{\theta}}{\mathrm{d} \theta} \right|_0 \implies -\left. \frac{p_0 q_1}{(1 + q_1 \theta)^2} \right|_0 = 1 \implies q_1 = -1  \]

Si ottiene quindi:

\[R_{(1,0)}(\theta) = \frac{1}{1 - \theta} \] 

%
\subsection{Approssimante $R_{(1,1)}(\theta)$}

Definiamo infine l’approssimante  $R_{(1,1)}(\theta)$  come:

$$R_{(1,1)}(\theta) = \frac{p_0 + p_1 \theta}{1 + q_1 \theta}$$  

Imponendo le condizioni di approssimazione:

$$R_{(1,1)}(0) = \left. e^\theta \right|_0 \implies p_0 = 1$$ 

\[\left. \frac{\mathrm{d} R_{(1,1)}}{\mathrm{d} \theta} \right|_0 = \left. \frac{\mathrm{d} e^\theta}{\mathrm{d} \theta} \right|_0 \implies \left. \frac{p_1 - q_1}{(1 + q_1 \theta)^2} \right|_0 = 1  \]
\[\left. \frac{\mathrm{d}^2 R_{(1,1)}}{\mathrm{d} \theta^2} \right|_0 = \left. \frac{\mathrm{d}^2 e^\theta}{\mathrm{d} \theta^2} \right|_0 \implies -\frac{2(1 + q_1 \theta) q_1 (p_1 - q_1)}{(1 + q_1 \theta)^4} \Bigg|_0 = 1  \]

Dalle equazioni precedenti si ricava il sistema:

\[\begin{aligned}
	p_1 - q_1 &= 1\\
	-2 q_1 (p_1 - q_1) &= 1
\end{aligned}\]  

La cui soluzione è: $p_1 = \frac{1}{2},\ q_1 = -\frac{1}{2}$ .

Si ottiene quindi:

\[R_{(1,1)}(\theta) = \frac{1 + \frac{1}{2} \theta}{1 - \frac{1}{2} \theta} \]
